\noindent \begin{center}
\section*{Resumen}
\par\end{center}

ROS proporciona un conjunto de herramientas y librerías para el desarrollo de aplicaciones robóticas 
que facilitan la comunicación entre diferentes componentes del sistema, así como la integración de sensores y actuadores.
Existen diversas aplicaciones que permiten la visualización y control de robots, sin embargo, la mayoría están centradas en entornos de escritorio.
Esto limita la capacidad de los usuarios para interactuar con los robots de manera remota y en tiempo real, especialmente en situaciones 
donde se requiere movilidad y flexibilidad.
En este trabajo se desarrolla una aplicación móvil multiplataforma de realidad aumentada para la visualización de variables de robots utilizando ROS2 y Unity,
de manera que los usuarios puedan acceder a la información del robot desde cualquier dispositivo, incluyendo smartphones y tabletas,
reduciendo la barrera entre el mundo físico y el digital y permitiendo una interacción más intuitiva y eficiente con los robots.

\vspace{0.5cm}
\begin{flushleft}
\textbf{Palabras clave}: teleoperación, robot, ROS2, aplicación, Unity, smartphone, 3D, ROS, multiplataforma, visualización, control, Android, iOS, realidad aumentada
\end{flushleft}